\documentclass[11pt]{article}

\usepackage[a4paper,margin=2.2cm]{geometry}
\usepackage{amsmath,amssymb}
\usepackage{booktabs}
\usepackage{graphicx}
\usepackage{multirow}
\usepackage{float}
\usepackage[hidelinks]{hyperref}
\usepackage{microtype}
\usepackage{url}

\graphicspath{{./}}
\renewcommand{\tablename}{Supplementary Table}
\renewcommand{\figurename}{Supplementary Figure}
\hypersetup{
  pdftitle={Supplementary Information: Sharp Target-Domain Certificates for Quantum-Kernel Advantage under Distribution Shift},
  pdfauthor={Roberto Fernandez-Barrios, Iker Pastor-Lopez, Asier Gonzalez-Santocildes, Pablo Garcia Bringas},
  pdfsubject={Supplementary Information (Additional file 1) submitted to EPJ Quantum Technology},
  pdfkeywords={quantum machine learning, quantum kernels, partial identification, active testing, distribution shift},
  hidelinks
}

\title{\textbf{Supplementary Information}\\[0.5em]
Sharp Target-Domain Certificates for Quantum-Kernel Advantage under Distribution Shift}
\author{Roberto Fern\'andez-Barrios, Iker Pastor-L\'opez,\\
Asier Gonz\'alez-Santocildes, and Pablo Garc\'ia Bringas\\[0.4em]
\small Faculty of Engineering, University of Deusto, Bilbao, Spain}
\date{}

\begin{document}
\maketitle

\section*{Supplementary Results: Scope}

This document contains extended results, complete sensitivity tables, and
diagnostics supporting the main article. No procedure is defined exclusively
here: all data construction, preprocessing, candidate grids, selection rules,
inference, external acquisition, geometry calculations, and finite-shot
procedures are reported in the main manuscript Methods. Machine-readable
values and frozen input hashes are included in the immutable release artifact.

\section*{Supplementary Results: Target-label evidence frontier}
\label{supp:v9_frontier}

The main text reports the 0.010 materiality threshold. Supplementary
Table~\ref{supp:tab_v9_thresholds} gives the fixed adaptive audit at 0.005 and
0.020; Supplementary Table~\ref{supp:tab_v9_acquisition} compares strong
prediction-aware controls, the exact ex-post oracle, and the weak
prediction-independent reference at 0.010. A value of zero means the
zero-label certificate already crosses the threshold. Every adaptive count is
deterministic for the fixed prediction matrix, labels, acquisition rule, and
tie order. Comparator medians summarize 200 prespecified SHA-256 tie orders
and are algorithmic sensitivities, not confidence intervals.

The prespecified reference-breadth sensitivity uses complete five-dimension
kernel blocks and 5,000 SHA-256-derived nested block orderings. It therefore
contains three evaluated breadths rather than only the customary and full
endpoints. The final Supplementary table summarizes the median zero-label
endpoint within each case--classifier cell over those orderings and then
across the 16 cells. Candidate count is a measure of reference breadth only,
not time, memory, or computational cost.

\begin{table}[H]
\centering
\caption{Target labels required against the full 115-candidate classical
family. Columns $n_{.005}$, $n_{.010}$, and $n_{.020}$ use the fixed adaptive
coverage audit. Every target batch contains 500 points.}
\label{supp:tab_v9_thresholds}
\scriptsize
\setlength{\tabcolsep}{4pt}
\begin{tabular}{llrrr}
\toprule
Target case & Classifier & $n_{.005}$ & $n_{.010}$ & $n_{.020}$\\
\midrule
EMBER $m1$ & SVC & 15 & 12 & 10\\
EMBER $m2$ & SVC & 17 & 15 & 12\\
ToN-IoT constructed & SVC & 0 & 0 & 0\\
ToN-IoT campaign & SVC & 0 & 0 & 0\\
UNSW DoS constructed & SVC & 3 & 0 & 0\\
UNSW DoS campaign & SVC & 13 & 11 & 5\\
UNSW Recon constructed & SVC & 17 & 14 & 11\\
UNSW Recon campaign & SVC & 34 & 28 & 11\\
\addlinespace
EMBER $m1$ & GPC & 32 & 29 & 26\\
EMBER $m2$ & GPC & 40 & 33 & 23\\
ToN-IoT constructed & GPC & 4 & 3 & 0\\
ToN-IoT campaign & GPC & 1 & 0 & 0\\
UNSW DoS constructed & GPC & 24 & 22 & 20\\
UNSW DoS campaign & GPC & 14 & 12 & 5\\
UNSW Recon constructed & GPC & 23 & 22 & 19\\
UNSW Recon campaign & GPC & 54 & 14 & 10\\
\bottomrule
\end{tabular}
\end{table}

\begin{table}[H]
\centering
\caption{Acquisition controls at $U\leq0.010$. Oracle is the exact
retrospective minimum; adaptive is the fixed bottleneck-coverage rule; random
active and initial coverage are medians across 200 SHA-256 tie orders; hash
all is the weak prediction-independent median. The adaptive rule does not
dominate either strong control.}
\label{supp:tab_v9_acquisition}
\scriptsize
\setlength{\tabcolsep}{3.2pt}
\begin{tabular}{llrrrrr}
\toprule
Target case & Classifier & Oracle & Adaptive & Random active & Initial coverage & Hash all\\
\midrule
EMBER $m1$ & SVC & 10 & 12 & 14.0 & 14.0 & 267.0\\
EMBER $m2$ & SVC & 11 & 15 & 16.0 & 16.0 & 284.0\\
ToN-IoT constructed & SVC & 0 & 0 & 0.0 & 0.0 & 0.0\\
ToN-IoT campaign & SVC & 0 & 0 & 0.0 & 0.0 & 0.0\\
UNSW DoS constructed & SVC & 0 & 0 & 0.0 & 0.0 & 0.0\\
UNSW DoS campaign & SVC & 6 & 11 & 8.0 & 8.0 & 220.0\\
UNSW Recon constructed & SVC & 9 & 14 & 13.0 & 13.0 & 234.0\\
UNSW Recon campaign & SVC & 7 & 28 & 28.0 & 286.5 & 312.5\\
\addlinespace
EMBER $m1$ & GPC & 18 & 29 & 34.0 & 34.0 & 369.5\\
EMBER $m2$ & GPC & 16 & 33 & 34.0 & 34.0 & 371.5\\
ToN-IoT constructed & GPC & 2 & 3 & 2.0 & 2.0 & 112.0\\
ToN-IoT campaign & GPC & 0 & 0 & 0.0 & 0.0 & 0.0\\
UNSW DoS constructed & GPC & 14 & 22 & 21.0 & 21.0 & 244.5\\
UNSW DoS campaign & GPC & 5 & 12 & 11.0 & 11.0 & 246.5\\
UNSW Recon constructed & GPC & 20 & 22 & 25.0 & 24.0 & 260.0\\
UNSW Recon campaign & GPC & 5 & 14 & 12.0 & 12.0 & 352.5\\
\midrule
Median over cells & & 6.5 & 13.0 & 12.5 & 12.5 & 245.5\\
\bottomrule
\end{tabular}
\end{table}

The frozen initial-coverage policy has one severe failure: for UNSW Recon
campaign SVC, its median is 286.5 labels because the zero-label witness set
becomes stale. Recomputing the actionable witness set avoids that failure,
but optimizing tied-witness coverage gives no consistent advantage over
dynamic random-active sampling.

\begin{figure}[H]
\centering
\includegraphics[width=\textwidth]{fig_v9_label_frontier.pdf}
\caption{Witness-directed contraction of the sharp accuracy upper endpoint
against the full classical family. Each curve is one fixed target batch and
selected quantum model; colour denotes source dataset, and solid/dashed lines
separate its paired shifts. Labels audit the fixed comparison and never
retune or select a model. Curves can cross below zero because the quantum
classifier can become certifiably worse than the best classical witness. At
500 labels every curve equals the realized advantage.}
\label{supp:fig_v9_frontier}
\end{figure}

\begin{table}[H]
\centering
\caption{Canonical fixed runs and P1$'$ quantum configurations used by every
target certificate. The exact run identifier is the displayed prefix followed
by \texttt{\_\_ms42\_\_q1000\_id500\_ood500\_\_qs42\_\_s42}; a slash in a
prefix denotes a literal \texttt{\_\_} separator. Configuration columns
identify the classifier and report map, scale, dimension, and the selected
SVC $C$ where applicable.}
\label{supp:tab_v9_runs}
\scriptsize
\setlength{\tabcolsep}{3.5pt}
\begin{tabular}{llll}
\toprule
Target case & Run prefix & SVC configuration & GPC configuration\\
\midrule
EMBER $m1$ & \texttt{m1\_hist\_byteent} & \texttt{pauli\_xz\_r1\_full/as2/d12/C10} & \texttt{zz\_r1\_full/as0.5/d12}\\
EMBER $m2$ & \texttt{m2\_hist\_byteent} & \texttt{pauli\_xz\_r1\_full/as2/d12/C100} & \texttt{zz\_r2\_full/as2/d8}\\
ToN-IoT constructed & \texttt{toniot\_scanning/m2\_centroid} & \texttt{zz\_r2\_full/as2/d6/C10} & \texttt{zz\_r1\_full/as2/d6}\\
ToN-IoT campaign & \texttt{toniot\_scanning/natural\_cur} & \texttt{pauli\_xz\_r1\_full/as2/d12/C10} & \texttt{zz\_r2\_full/as0.5/d10}\\
UNSW DoS constructed & \texttt{unsw\_dos/m2\_centroid} & \texttt{zmap\_r2/as0.5/d12/C100} & \texttt{zz\_r1\_full/d12}\\
UNSW DoS campaign & \texttt{unsw\_dos/natural\_cur} & \texttt{zz\_r2\_full/as0.5/d4/C100} & \texttt{zz\_r1\_full/as0.5/d4}\\
UNSW Recon constructed & \texttt{unsw\_recon/m2\_centroid} & \texttt{zz\_r2\_full/as0.5/d4/C100} & \texttt{zz\_r2\_full/as0.5/d8}\\
UNSW Recon campaign & \texttt{unsw\_recon/natural\_cur} & \texttt{zz\_r1\_full/as0.5/d6/C10} & \texttt{zz\_r2\_full/as0.5/d4}\\
\bottomrule
\end{tabular}
\end{table}

\section*{Supplementary Results: Prevalence-conditional balanced accuracy}
\label{supp:v9_bacc}

Supplementary Table~\ref{supp:tab_v9_bacc} reports the exact finite-batch
MILP after fixing prevalence to the target value that becomes available when
the label file is unlocked. It is a retrospective metric sensitivity, not a
label-free use of the true prevalence. The LP relaxation and integer program
agree to numerical precision in all reported cells. Because these batches
are close to balanced by construction, the balanced-accuracy upper endpoints
closely track their accuracy counterparts and do not rescue a broad
zero-label certificate.

\begin{table}[htbp]
\centering
\caption{Prevalence-conditional sharp balanced-accuracy upper endpoints
against the full classical family. $U_{\rm Acc}$ is the assumption-free
zero-label accuracy endpoint, $U_{\rm BAcc}(\pi)$ is the exact MILP endpoint
at the retrospectively unlocked prevalence, and
$\Delta_{\rm BAcc,full}$ is the realized balanced-accuracy advantage.}
\label{supp:tab_v9_bacc}
\scriptsize
\setlength{\tabcolsep}{4pt}
\begin{tabular}{llrrrr}
\toprule
Target case & Classifier & $\pi$ & $U_{\rm Acc}$ & $U_{\rm BAcc}(\pi)$ &
$\Delta_{\rm BAcc,full}$\\
\midrule
EMBER $m1$ & SVC & .488 & .050 & .050 & $-.024$\\
EMBER $m2$ & SVC & .488 & .054 & .054 & $-.037$\\
ToN-IoT constructed & SVC & .492 & .002 & .002 & $-.004$\\
ToN-IoT campaign & SVC & .518 & .002 & .002 & $-.018$\\
UNSW DoS constructed & SVC & .496 & .010 & .010 & $-.016$\\
UNSW DoS campaign & SVC & .520 & .034 & .035 & $-.025$\\
UNSW Recon constructed & SVC & .498 & .046 & .046 & $-.032$\\
UNSW Recon campaign & SVC & .470 & .034 & .033 & $-.022$\\
\addlinespace
EMBER $m1$ & GPC & .488 & .080 & .078 & $-.033$\\
EMBER $m2$ & GPC & .488 & .074 & .073 & $-.016$\\
ToN-IoT constructed & GPC & .492 & .016 & .016 & $-.014$\\
ToN-IoT campaign & GPC & .518 & .006 & .006 & $-.012$\\
UNSW DoS constructed & GPC & .496 & .066 & .066 & $-.052$\\
UNSW DoS campaign & GPC & .520 & .028 & .029 & $-.033$\\
UNSW Recon constructed & GPC & .498 & .088 & .088 & $-.066$\\
UNSW Recon campaign & GPC & .470 & .030 & .029 & $+.001$\\
\bottomrule
\end{tabular}
\end{table}

\section*{Supplementary Results: Prospective Gate-2 replication against the prespecified classical-kernel reference family}
\label{supp:v10_gate2}

The Gate-2 protocol against the prespecified classical-kernel reference family was fixed before acquiring any of its three tasks,
computing a model prediction, or opening a target label. Its specification
SHA-256 is
\texttt{3a8318d92d4af2ae\allowbreak eaf0c0edb069c3be\allowbreak 59f31da6d1ee50fb\allowbreak 6a6256e9d9d280b0}.
The prediction stage copied target labels into a separate sealed tree and
then removed them from the fitting state. Once every available unit contained
175 candidate checkpoints for both classifiers, the aggregate prediction
lock was written with SHA-256
\texttt{e4ee6bb89bc78047\allowbreak e7c848a0db3a6255\allowbreak cf0a2600a21b42fb\allowbreak ab956a0e889507d1}.
Only then did the audit process record its one-time label opening. The tasks
were inherited unchanged from the unused fallback list in the earlier
external protocol.

During acquisition, the pinned NHANES download URLs returned HTML 404 pages.
The transport paths alone were updated to the current official CDC directory;
the same 37 cycle-specific XPT filenames, TableShift parser, columns, and
domain definition were retained. XPORT magic bytes and SHA-256 hashes were
verified before export, and the complete URL-repair audit is archived. This
maintenance preceded every model execution and did not alter the subsequent
minimum-feature disposition.

\begin{table}[H]
\centering
\caption{Technical disposition of every prospective Gate-2 task against the prespecified classical-kernel reference family.
``Candidates'' counts kernel configurations per available seed; each
configuration supplies SVC and GPC predictions. NHANES was retained as a
recorded technical failure and was not replaced.}
\label{supp:tab_v10_availability}
\scriptsize
\setlength{\tabcolsep}{4pt}
\begin{tabular}{lclrr}
\toprule
Task & Shift & Status & Seeds & Candidates/seed\\
\midrule
BRFSS Diabetes & race & prediction-locked & 5 & 175\\
ACS Food Stamps & Census division & prediction-locked & 5 & 175\\
NHANES Lead & poverty & minimum-feature failure ($11<12$) & 5 & 0\\
\bottomrule
\end{tabular}
\end{table}

Supplementary Table~\ref{supp:tab_v10_primary} gives every seed-level primary
cell. The sharp endpoint is already at most 0.010 without labels in 11/20
cells. Across all cells, the adaptive audit has median zero and range 0--76;
the exact label oracle has median zero and range 0--40. Every realized
accuracy effect is negative. The one broad zero-label case, ACS Food Stamps
SVC at seed 999, is retained and determines the maximum adaptive count.

\begin{table}[H]
\centering
\caption{Every prospective full-family cell at the primary accuracy threshold.
$U_0$ is the zero-label sharp upper endpoint, $n_{.01}$ is the fixed adaptive
count, $n^*_{.01}$ is the exact retrospective label-oracle minimum, and
$\Delta_{\rm full}$ is the realized quantum-minus-best-classical accuracy
difference after label opening.}
\label{supp:tab_v10_primary}
\scriptsize
\setlength{\tabcolsep}{3.5pt}
\begin{tabular}{llrrrrr}
\toprule
Task & Classifier & Seed & $U_0$ & $n_{.01}$ & $n^*_{.01}$ & $\Delta_{\rm full}$\\
\midrule
BRFSS Diabetes & SVC & 42   & .012 &  3 &  1 & $-.154$\\
                 & SVC & 123  & .008 &  0 &  0 & $-.156$\\
                 & SVC & 999  & .008 &  0 &  0 & $-.110$\\
                 & SVC & 7    & .024 & 12 &  4 & $-.096$\\
                 & SVC & 2024 & .018 &  4 &  2 & $-.118$\\
                 & GPC & 42   & .000 &  0 &  0 & $-.008$\\
                 & GPC & 123  & .002 &  0 &  0 & $-.010$\\
                 & GPC & 999  & .004 &  0 &  0 & $-.018$\\
                 & GPC & 7    & .002 &  0 &  0 & $-.002$\\
                 & GPC & 2024 & .018 &  2 &  2 & $-.026$\\
\addlinespace
ACS Food Stamps & SVC & 42   & .010 &  0 &  0 & $-.030$\\
                 & SVC & 123  & .014 &  1 &  1 & $-.060$\\
                 & SVC & 999  & .168 & 76 & 40 & $-.104$\\
                 & SVC & 7    & .010 &  0 &  0 & $-.108$\\
                 & SVC & 2024 & .006 &  0 &  0 & $-.076$\\
                 & GPC & 42   & .010 &  0 &  0 & $-.004$\\
                 & GPC & 123  & .006 &  0 &  0 & $-.018$\\
                 & GPC & 999  & .022 &  3 &  3 & $-.014$\\
                 & GPC & 7    & .032 & 10 &  6 & $-.010$\\
                 & GPC & 2024 & .024 &  4 &  4 & $-.006$\\
\bottomrule
\end{tabular}
\end{table}

The predefined outcome summaries are the medians over five seeds in each
task--classifier cell. They are 3, 0, 0, and 3 labels for BRFSS SVC, BRFSS
GPC, ACS SVC, and ACS GPC, respectively. All four are at most 50, the overall
seed-level median is zero, and the maximum is 76, satisfying every numerical
condition in the predefined primary criterion. Because only two tasks
completed the grid, this is reported factually as prospective corroboration
in two technically eligible tasks, not as a three-task result.

\begin{table}[H]
\centering
\caption{Prospective acquisition controls at $U\leq0.010$. Adaptive entries
are median [minimum, maximum] over five seeds. Other entries are median
[2.5th, 97.5th percentile] over 1,000 seed--order combinations. Ranges are
algorithmic sensitivities, not confidence intervals.}
\label{supp:tab_v10_controls}
\scriptsize
\setlength{\tabcolsep}{3.4pt}
\begin{tabular}{llrrrr}
\toprule
Task & Classifier & Adaptive & Random active & Initial coverage & Hash all\\
\midrule
BRFSS Diabetes & SVC & $3[0,12]$ & $2[0,12]$ & $2[0,12]$ & $91[0,301]$\\
                & GPC & $0[0,2]$  & $0[0,2]$  & $0[0,2]$  & $0[0,162]$\\
ACS Food Stamps & SVC & $0[0,76]$ & $0[0,79]$ & $0[0,78]$ & $0[0,350]$\\
                & GPC & $3[0,10]$ & $4[0,15]$ & $4[0,15]$ & $141[0,415]$\\
\bottomrule
\end{tabular}
\end{table}

Adaptive coverage wins, ties, and loses against each strong comparator in 3,
15, and 2 seed-level cells; its median paired difference is zero. This
prospectively repeats the conclusion that actionable-disagreement restriction
is useful while coverage optimization is not demonstrably superior. The
realized-prevalence balanced-accuracy MILP is retained as a retrospective
sensitivity. Its full-family upper endpoints span 0--0.223, whereas every
realized balanced-accuracy effect is negative (range $-0.080$ to $-0.013$).
Because prevalence was not known
before label opening, these values are not prospective label-free
balanced-accuracy certificates.

All 20 selected prospective quantum models are product maps. Consequently,
this experiment corroborates Gate~2 against the prespecified classical-kernel
reference family on the eligible fixed tasks but
does not independently replicate the retrospective entangling-ZZ result.

\section*{Supplementary Results: Cross-protocol sensitivity map}
\label{supp:specification}

The ten rows remain in their prespecified S1--S10 order rather than being
ranked by outcome. Every row includes all eight fixed security
scenario-groups. Scenario-groups are averaged within EMBER, UNSW-NB15, and
ToN-IoT before the three source datasets receive equal weight. The corrected
S10 values in Supplementary Table~\ref{supp:tab_specification} are identical to the
confirmatory endpoint in the main article. This display compares complete
analysis specifications across the earlier fixed-$C$ protocol and the controlled protocol; it is not a factorial
decomposition of their individual choices.

\begin{table}[ht]
\centering
\caption{Prespecified cross-protocol definitions and source-dataset-equal
quantum-minus-classical OOD balanced-accuracy effects. ``Oracle'' selects and
evaluates on the same OOD test labels; ``ID-val'' uses a disjoint
in-distribution validation split. SVC regularization is fixed at $C=1$ or
selected by training-only cross-validation; this axis is not applicable to
GPC. ``ID-test'' selects on the in-distribution test split of the earlier fixed-$C$ runs, which had no disjoint validation half. S1--S4 use that earlier protocol and S5--S10 use the controlled protocol, so their boundary
does not isolate the effect of regularization. Values are descriptive
fixed-case summaries, not per-axis causal effects.}
\label{supp:tab_specification}
\scriptsize
\setlength{\tabcolsep}{3pt}
\begin{tabular}{lllllrr}
\toprule
ID & Selection & SVC $C$ & Classical reference & Budget & SVC & GPC \\
\midrule
S1  & OOD oracle & Fixed & Linear/RBF & Native & $+0.0313$ & $+0.0340$ \\
S2  & ID-test & Fixed & Linear/RBF & Native & $+0.0475$ & $+0.0418$ \\
S3  & OOD oracle & Fixed & Extended & Native & $+0.0058$ & $+0.0145$ \\
S4  & ID-test & Fixed & Extended & Native & $+0.0014$ & $+0.0122$ \\
S5  & OOD oracle & Train CV & Linear/RBF & Native & $+0.0034$ & $+0.0125$ \\
S6  & ID-val & Train CV & Linear/RBF & Native & $-0.0041$ & $+0.0059$ \\
S7  & OOD oracle & Train CV & Extended & Native & $-0.0054$ & $+0.0012$ \\
S8  & ID-val & Train CV & Extended & Native & $-0.0038$ & $-0.0029$ \\
S9  & OOD oracle & Train CV & Extended & Equal & $-0.0032$ & $+0.0038$ \\
S10 & ID-val & Train CV & Extended & Equal & $-0.0056$ & $-0.0009$ \\
\bottomrule
\end{tabular}
\end{table}

The source-dataset-equal range is 0.0531 for SVC and 0.0447 for GPC.
This spread is not an interval: the specifications share data and candidates
and deliberately target different estimands.

\section*{Supplementary Results: Candidate-budget sampling}
\label{supp:budget}

Supplementary Table~\ref{supp:tab_budget} shows every group--classifier cell under all three
prespecified sampling schemes. The primary blocked scheme mirrors the
five-dimension block structure of the quantum pool; the other schemes match
candidate count but not the same search-space structure. Their close values
show that the controlled conclusion is not an artefact of the blocked draw.

\begin{table}[ht]
\centering
\caption{Complete equal-budget sampling sensitivity. Entries are median
quantum-minus-classical P1$'$ OOD differences over 5,000 classical-pool
subsamples. ``Blocked'' is primary; ``uniform'' and ``stratified'' are
sensitivities. Source-dataset-equal rows average scenario-group medians within
the three source datasets. These resampling medians are distinct from the
run-level expected-classical endpoint used for the conditional intervals in
the main article.}
\label{supp:tab_budget}
\small
\setlength{\tabcolsep}{4.5pt}
\begin{tabular}{llrrr}
\toprule
Scenario & Clf. & Blocked & Uniform & Stratified \\
\midrule
\multirow{2}{*}{EMBER $m1$}
 & SVC & $-0.0056$ & $-0.0054$ & $-0.0054$ \\
 & GPC & $-0.0004$ & $-0.0009$ & $-0.0013$ \\
\multirow{2}{*}{EMBER $m2$}
 & SVC & $-0.0071$ & $-0.0071$ & $-0.0071$ \\
 & GPC & $+0.0007$ & $+0.0008$ & $+0.0005$ \\
\multirow{2}{*}{UNSW-DoS (campaign)}
 & SVC & $-0.0036$ & $-0.0034$ & $-0.0036$ \\
 & GPC & $-0.0072$ & $-0.0072$ & $-0.0071$ \\
\multirow{2}{*}{UNSW-DoS (constructed)}
 & SVC & $-0.0046$ & $-0.0030$ & $-0.0028$ \\
 & GPC & $-0.0116$ & $-0.0113$ & $-0.0112$ \\
\multirow{2}{*}{UNSW-Recon (campaign)}
 & SVC & $+0.0050$ & $+0.0052$ & $+0.0052$ \\
 & GPC & $+0.0046$ & $+0.0047$ & $+0.0049$ \\
\multirow{2}{*}{UNSW-Recon (constructed)}
 & SVC & $-0.0090$ & $-0.0089$ & $-0.0089$ \\
 & GPC & $-0.0082$ & $-0.0076$ & $-0.0077$ \\
\multirow{2}{*}{ToN-IoT (campaign)}
 & SVC & $-0.0147$ & $-0.0149$ & $-0.0143$ \\
 & GPC & $-0.0095$ & $-0.0094$ & $-0.0096$ \\
\multirow{2}{*}{ToN-IoT (constructed)}
 & SVC & $-0.0011$ & $-0.0008$ & $-0.0011$ \\
 & GPC & $+0.0093$ & $+0.0063$ & $+0.0057$ \\
\midrule
\multirow{2}{*}{Source-dataset-equal}
 & SVC & $-0.00576$ & $-0.00555$ & $-0.00548$ \\
 & GPC & $-0.00185$ & $-0.00232$ & $-0.00253$ \\
\bottomrule
\end{tabular}
\end{table}

The largest within-cell scheme range is 0.00353 for ToN-IoT constructed
under GPC. The next largest is 0.00184 for UNSW-DoS constructed under SVC.

\section*{Supplementary Results: Effective-rank matching sensitivities}
\label{supp:rank}

The predefined rank ratio is
$\max(r_Q,r_C)/\min(r_Q,r_C)$. Supplementary Table~\ref{supp:tab_rank} reports all
calipers for both observational matchers. With replacement, 75.2\% of pairs
survive the primary 1.25 caliper. One-to-one assignment has a smaller common
support and becomes visibly poor without a caliper (95th-percentile ratio
11.45), making the overlap restriction scientifically material.

\begin{table}[ht]
\centering
\caption{Pooled SVC effective-rank matching sensitivity. ``Nearest'' pairs
each quantum candidate to its nearest classical rank with classical reuse;
``one-to-one'' uses minimum-cost assignment without reuse. Retained
percentages use 64,800 possible quantum candidates as denominator.
$\Delta$ is quantum minus classical OOD balanced accuracy.}
\label{supp:tab_rank}
\scriptsize
\setlength{\tabcolsep}{3.5pt}
\begin{tabular}{llrrrrr}
\toprule
Matcher & Caliper & Pairs (\%) & Median ratio & 95th pct.\ ratio &
Median $\Delta$ & Q better \\
\midrule
\multirow{5}{*}{Nearest}
 & 1.10 & 27,928 (43.1) & 1.040 & 1.093 & $-0.0054$ & 29.4\% \\
 & 1.25 & 48,708 (75.2) & 1.082 & 1.226 & $-0.0056$ & 29.7\% \\
 & 1.50 & 60,647 (93.6) & 1.112 & 1.390 & $-0.0061$ & 28.9\% \\
 & 2.00 & 64,408 (99.4) & 1.123 & 1.542 & $-0.0070$ & 27.8\% \\
 & none & 64,800 (100) & 1.124 & 1.578 & $-0.0071$ & 27.6\% \\
\midrule
\multirow{5}{*}{One-to-one}
 & 1.10 & 13,205 (20.4) & 1.044 & 1.094 & $-0.0048$ & 30.2\% \\
 & 1.25 & 24,675 (38.1) & 1.092 & 1.232 & $-0.0053$ & 30.3\% \\
 & 1.50 & 35,129 (54.2) & 1.151 & 1.447 & $-0.0057$ & 30.2\% \\
 & 2.00 & 46,898 (72.4) & 1.232 & 1.829 & $-0.0057$ & 30.8\% \\
 & none & 64,800 (100) & 1.418 & 11.452 & $-0.0062$ & 30.9\% \\
\bottomrule
\end{tabular}
\end{table}

At the primary 1.25 caliper, both matchers yield a negative group median in
all eight groups. The result is descriptive of the retained overlap, not a
causal family estimate.

\section*{Supplementary Results: Geometry associations}
\label{supp:geometry}

Effective-rank association changes sign across regimes, whereas OOD
kernel-target alignment is more consistently positive
(Supplementary Table~\ref{supp:tab_geometry}). The source-dataset leave-one-out ordering
correlations are 0.28 for EMBER, 0.22 for UNSW-NB15, and 0.40 for ToN-IoT
when trained on the other two sources.

\begin{table}[ht]
\centering
\caption{Geometry--OOD association and primary nearest-rank pairing for SVC.
``Partial'' controls kernel family and scale. Matching uses replacement and
rank ratio $\leq1.25$; $\Delta$ is the median quantum-minus-classical OOD
difference. All quantities are observational.}
\label{supp:tab_geometry}
\small
\setlength{\tabcolsep}{4pt}
\begin{tabular}{lccccc}
\toprule
Scenario & $\rho$ rank & Partial & $\rho$ KTA & Matched $\Delta$ & Q better \\
\midrule
EMBER $m1$ & $+0.32$ & $+0.36$ & $+0.30$ & $-0.0094$ & 33\% \\
EMBER $m2$ & $+0.80$ & $+0.82$ & $+0.63$ & $-0.0089$ & 22\% \\
UNSW-DoS (campaign) & $+0.44$ & $+0.41$ & $+0.64$ & $-0.0089$ & 27\% \\
UNSW-DoS (constructed) & $-0.14$ & $-0.08$ & $+0.39$ & $-0.0057$ & 29\% \\
UNSW-Recon (campaign) & $+0.06$ & $+0.06$ & $+0.22$ & $-0.0047$ & 37\% \\
UNSW-Recon (constructed) & $-0.15$ & $-0.08$ & $+0.37$ & $-0.0061$ & 30\% \\
ToN-IoT (campaign) & $+0.25$ & $+0.27$ & $+0.46$ & $-0.0060$ & 24\% \\
ToN-IoT (constructed) & $+0.65$ & $+0.63$ & $+0.59$ & $-0.0013$ & 35\% \\
\bottomrule
\end{tabular}
\end{table}

\section*{Supplementary Results: Descriptive oracle comparison}
\label{supp:oracle}

The initial EMBER evaluation selected its representative using OOD test
labels and fixed SVC at $C=1$. Expanding the reference family nearly removes
the apparent gain even before target-label-free selection is imposed.

\begin{table}[ht]
\centering
\caption{Descriptive EMBER oracle results under fixed $C=1$ (18 correlated
settings). Mean and median effects are quantum minus classical balanced
accuracy. These values are not deployable estimates or independent
replicates.}
\label{supp:tab_oracle}
\small
\begin{tabular}{llcccc}
\toprule
Selection & Classifier & Classical reference & Q wins & Mean $\Delta$ & Median $\Delta$ \\
\midrule
\multirow{4}{*}{Best-by-OOD}
 & SVC & linear+RBF & 17/18 & $+0.0374$ & $+0.0426$ \\
 & SVC & extended & 15/18 & $+0.0044$ & $+0.0064$ \\
 & GPC & linear+RBF & 15/18 & $+0.0283$ & $+0.0339$ \\
 & GPC & extended & 11/18 & $+0.0023$ & $+0.0029$ \\
\midrule
\multirow{4}{*}{Best-by-ID}
 & SVC & linear+RBF & 18/18 & $+0.0775$ & $+0.0794$ \\
 & SVC & extended & 18/18 & $+0.0149$ & $+0.0152$ \\
 & GPC & linear+RBF & 18/18 & $+0.0649$ & $+0.0664$ \\
 & GPC & extended & 18/18 & $+0.0192$ & $+0.0166$ \\
\bottomrule
\end{tabular}
\end{table}

\section*{Supplementary Results: Small-cluster interval sensitivity}
\label{supp:intervals}

Supplementary Table~\ref{supp:tab_intervals} compares the primary pointwise conditional
cluster-$t$ interval with the prespecified BCa and leave-one-q-split
diagnostics. The effective replication unit is five q-split clusters in every
row. These intervals are not simultaneous and are not used as a
familywise-error-controlled classification of the eight cases. Differences
between interval constructions do not change the qualitative fixed-case
description, and an interval containing zero is not interpreted as
equivalence. With only five clusters, the BCa interval is retained as a
fragile sensitivity rather than treated as additional confirmatory evidence.

\begin{table}[ht]
\centering
\caption{Pointwise conditional uncertainty sensitivity for all primary equal-budget
group--classifier cells. $\Delta$ is quantum minus classical OOD balanced
accuracy. The cluster-$t$ and 9,999-draw BCa columns are intervals over the
same five q-split cluster means ($n=5$); LOO is the range after leaving out
one q-split cluster. All are conditional on the fixed benchmark pools and
design, not intervals over a dataset population. They are not simultaneous
intervals, and the BCa column is fragile at this effective sample size.}
\label{supp:tab_intervals}
\scriptsize
\setlength{\tabcolsep}{3.0pt}
\begin{tabular}{llrrrr}
\toprule
Scenario & Clf. & $\Delta$ & 95\% cluster-$t$ & 95\% BCa & LOO range \\
\midrule
\multirow{2}{*}{EMBER $m1$}
 & GPC & $-0.0009$ & $[-0.0101,+0.0083]$ & $[-0.0062,+0.0054]$ & $[-0.0035,+0.0015]$ \\
 & SVC & $-0.0058$ & $[-0.0340,+0.0223]$ & $[-0.0226,+0.0110]$ & $[-0.0120,+0.0014]$ \\
\multirow{2}{*}{EMBER $m2$}
 & GPC & $+0.0027$ & $[-0.0027,+0.0081]$ & $[-0.0015,+0.0053]$ & $[+0.0017,+0.0043]$ \\
 & SVC & $-0.0064$ & $[-0.0137,+0.0009]$ & $[-0.0107,-0.0018]$ & $[-0.0081,-0.0049]$ \\
\multirow{2}{*}{UNSW-DoS (campaign)}
 & GPC & $-0.0057$ & $[-0.0120,+0.0005]$ & $[-0.0084,+0.0004]$ & $[-0.0078,-0.0046]$ \\
 & SVC & $-0.0033$ & $[-0.0139,+0.0073]$ & $[-0.0101,+0.0029]$ & $[-0.0058,-0.0008]$ \\
\multirow{2}{*}{UNSW-DoS (constructed)}
 & GPC & $-0.0108$ & $[-0.0177,-0.0039]$ & $[-0.0141,-0.0054]$ & $[-0.0131,-0.0095]$ \\
 & SVC & $-0.0044$ & $[-0.0119,+0.0031]$ & $[-0.0087,+0.0005]$ & $[-0.0067,-0.0026]$ \\
\multirow{2}{*}{UNSW-Recon (campaign)}
 & GPC & $+0.0052$ & $[-0.0019,+0.0123]$ & $[-0.0001,+0.0089]$ & $[+0.0038,+0.0072]$ \\
 & SVC & $+0.0050$ & $[-0.0058,+0.0157]$ & $[+0.0006,+0.0131]$ & $[+0.0011,+0.0065]$ \\
\multirow{2}{*}{UNSW-Recon (constructed)}
 & GPC & $-0.0075$ & $[-0.0208,+0.0059]$ & $[-0.0162,+0.0006]$ & $[-0.0106,-0.0040]$ \\
 & SVC & $-0.0089$ & $[-0.0146,-0.0031]$ & $[-0.0123,-0.0053]$ & $[-0.0102,-0.0076]$ \\
\multirow{2}{*}{ToN-IoT (campaign)}
 & GPC & $-0.0090$ & $[-0.0190,+0.0010]$ & $[-0.0166,-0.0034]$ & $[-0.0106,-0.0060]$ \\
 & SVC & $-0.0148$ & $[-0.0268,-0.0029]$ & $[-0.0218,-0.0068]$ & $[-0.0180,-0.0129]$ \\
\multirow{2}{*}{ToN-IoT (constructed)}
 & GPC & $+0.0109$ & $[+0.0017,+0.0202]$ & $[+0.0060,+0.0178]$ & $[+0.0082,+0.0129]$ \\
 & SVC & $-0.0011$ & $[-0.0049,+0.0028]$ & $[-0.0033,+0.0013]$ & $[-0.0019,-0.0003]$ \\
\bottomrule
\end{tabular}
\end{table}

\begin{table}[p]
\centering
\caption{Five q-split cluster means underlying each primary pointwise
conditional interval. Values are quantum-minus-expected-classical OOD
balanced-accuracy differences under P1$'$ and 60-candidate blocked matching.
Lower-level model/master-seed and nested-size realizations are averaged within
each displayed cluster; the five columns, not the individual pipeline
realizations, are the effective replication units.}
\label{supp:tab_cluster_values}
\scriptsize
\setlength{\tabcolsep}{3.4pt}
\begin{tabular}{llrrrrr}
\toprule
Scenario & Clf. & $q=7$ & $q=42$ & $q=123$ & $q=999$ & $q=2024$ \\
\midrule
\multirow{2}{*}{EMBER $m1$}
 & SVC & $-0.01393$ & $+0.01886$ & $-0.03483$ & $-0.01476$ & $+0.01558$ \\
 & GPC & $-0.00320$ & $-0.00166$ & $+0.00135$ & $+0.00965$ & $-0.01066$ \\
\multirow{2}{*}{EMBER $m2$}
 & SVC & $-0.01221$ & $-0.00951$ & $-0.00064$ & $+0.00052$ & $-0.01010$ \\
 & GPC & $+0.00351$ & $+0.00663$ & $+0.00625$ & $+0.00098$ & $-0.00395$ \\
\multirow{2}{*}{UNSW-DoS (campaign)}
 & SVC & $-0.00005$ & $-0.01302$ & $+0.00666$ & $+0.00135$ & $-0.01134$ \\
 & GPC & $-0.00595$ & $-0.01046$ & $-0.00637$ & $-0.00859$ & $+0.00270$ \\
\multirow{2}{*}{UNSW-DoS (constructed)}
 & SVC & $-0.00466$ & $-0.00689$ & $-0.01185$ & $+0.00469$ & $-0.00340$ \\
 & GPC & $-0.01167$ & $-0.01352$ & $-0.01614$ & $-0.00153$ & $-0.01109$ \\
\multirow{2}{*}{UNSW-Recon (campaign)}
 & SVC & $+0.00246$ & $+0.00188$ & $-0.00109$ & $+0.02025$ & $+0.00127$ \\
 & GPC & $+0.00998$ & $+0.00553$ & $-0.00281$ & $+0.01107$ & $+0.00233$ \\
\multirow{2}{*}{UNSW-Recon (constructed)}
 & SVC & $-0.00424$ & $-0.00367$ & $-0.01388$ & $-0.01128$ & $-0.01122$ \\
 & GPC & $-0.01376$ & $+0.00122$ & $-0.00861$ & $-0.02125$ & $+0.00511$ \\
\multirow{2}{*}{ToN-IoT (campaign)}
 & SVC & $-0.02155$ & $-0.00193$ & $-0.00700$ & $-0.02105$ & $-0.02260$ \\
 & GPC & $-0.02081$ & $-0.00256$ & $-0.00313$ & $-0.01385$ & $-0.00446$ \\
\multirow{2}{*}{ToN-IoT (constructed)}
 & SVC & $+0.00217$ & $-0.00268$ & $-0.00324$ & $+0.00247$ & $-0.00403$ \\
 & GPC & $+0.02184$ & $+0.01485$ & $+0.00872$ & $+0.00624$ & $+0.00304$ \\
\bottomrule
\end{tabular}
\end{table}

\clearpage
\section*{Supplementary Results: Quantum-map strata}
\label{supp:quantum_strata}

Supplementary Table~\ref{supp:tab_quantum_strata} reports every fixed-case SVC effect after
separating the 30 entangling-ZZ and 30 separable-product candidates. Each
stratum is compared with the expected P1$'$ winner over 5,000 blocked draws
of six complete extended-classical blocks, also 30 candidates. The full pool
uses 60 candidates per family. Intervals are pointwise and conditional on five
q-split clusters; their inclusion of zero is not an equivalence decision.

\begin{table}[ht]
\centering
\caption{Circuit-stratified quantum-minus-classical OOD balanced-accuracy
effects under train-CV and P1$'$. Brackets for scenario-groups are pointwise
95\% conditional cluster-$t$ intervals ($n=5$ q-split clusters), not
simultaneous intervals. The source-dataset-equal row gives the mean followed
by the range of its three constituent source-dataset means, not a confidence
interval.}
\label{supp:tab_quantum_strata}
\scriptsize
\setlength{\tabcolsep}{3pt}
\begin{tabular}{lccc}
\toprule
Scenario & All maps ($B=60$) & Entangling ZZ ($B=30$) & Product maps ($B=30$) \\
\midrule
EMBER $m1$ & $-0.0057\;[-0.0340,+0.0225]$ & $-0.0063\;[-0.0279,+0.0153]$ & $+0.0036\;[-0.0132,+0.0204]$ \\
EMBER $m2$ & $-0.0065\;[-0.0138,+0.0008]$ & $-0.0053\;[-0.0130,+0.0024]$ & $-0.0135\;[-0.0166,-0.0104]$ \\
UNSW-DoS (campaign) & $-0.0033\;[-0.0139,+0.0073]$ & $-0.0129\;[-0.0184,-0.0073]$ & $-0.0016\;[-0.0116,+0.0083]$ \\
UNSW-DoS (constructed) & $-0.0045\;[-0.0119,+0.0030]$ & $-0.0149\;[-0.0299,+0.0001]$ & $-0.0007\;[-0.0063,+0.0049]$ \\
UNSW-Recon (campaign) & $+0.0049\;[-0.0058,+0.0157]$ & $-0.0014\;[-0.0070,+0.0041]$ & $+0.0057\;[-0.0023,+0.0137]$ \\
UNSW-Recon (constructed) & $-0.0088\;[-0.0145,-0.0032]$ & $-0.0171\;[-0.0286,-0.0056]$ & $-0.0025\;[-0.0066,+0.0015]$ \\
ToN-IoT (campaign) & $-0.0149\;[-0.0270,-0.0028]$ & $-0.0196\;[-0.0303,-0.0089]$ & $-0.0225\;[-0.0347,-0.0102]$ \\
ToN-IoT (constructed) & $-0.0011\;[-0.0050,+0.0028]$ & $-0.0013\;[-0.0047,+0.0021]$ & $-0.0009\;[-0.0036,+0.0018]$ \\
\midrule
Source-dataset-equal & $-0.0057\;[-0.0080,-0.0029]$ & $-0.0093\;[-0.0116,-0.0058]$ & $-0.0055\;[-0.0117,+0.0002]$ \\
\bottomrule
\end{tabular}
\end{table}

For P1$'$ selection, product maps account for 594/1,080 SVC winners and
348/1,080 GPC winners. Under same-test OOD selection, the corresponding
counts are 650/1,080 and 528/1,080. Complete counts by individual map are
provided in \texttt{quantum\_winner\_composition.csv} and shown in
Supplementary Table~\ref{supp:tab_quantum_winners}.

\begin{table}[htbp]
\centering
\caption{Quantum-map composition of the within-family winner. Each
model--selector column contains 1,080 fixed security runs. Counts and
percentages describe which configuration maximizes the specified selection
metric within the quantum family; they do not compare its OOD accuracy with
the classical family. P1$'$ uses disjoint ID validation, whereas the OOD
oracle reuses the shifted-test labels.}
\label{supp:tab_quantum_winners}
\scriptsize
\setlength{\tabcolsep}{3.2pt}
\begin{tabular}{llrrrr}
\toprule
Map & Stratum & P1$'$ SVC & P1$'$ GPC & OOD-oracle SVC & OOD-oracle GPC \\
\midrule
ZZ, 1 rep. & entangling ZZ & 290 (26.9\%) & 353 (32.7\%) & 282 (26.1\%) & 279 (25.8\%) \\
ZZ, 2 reps. & entangling ZZ & 196 (18.1\%) & 379 (35.1\%) & 148 (13.7\%) & 273 (25.3\%) \\
Pauli-X/Z, 1 rep. & separable product & 315 (29.2\%) & 171 (15.8\%) & 353 (32.7\%) & 246 (22.8\%) \\
Z, 2 reps. & separable product & 279 (25.8\%) & 177 (16.4\%) & 297 (27.5\%) & 282 (26.1\%) \\
\bottomrule
\end{tabular}
\end{table}

\section*{Supplementary Results: Evaluation-choice factorial}
\label{supp:factorial}

The q1000 factorial holds the controlled-protocol data generation, split roles, kernel
inventory, and preprocessing fixed while crossing regularization, selector,
classical reference, and budget mode. Supplementary Table~\ref{supp:tab_factorial} reports
all 16 three-source-dataset-equal endpoints. The fixed-$C$, same-test-oracle,
customary-reference, native-budget corner is $+0.0131$; the train-CV,
ID-validation, extended-reference, equal-count corner is $-0.0047$.
These endpoints are fixed-case summaries, not population effects.

\begin{table}[htbp]
\centering
\caption{Evaluation-choice q1000 SVC factorial. Entries are three-source-dataset-equal
quantum-minus-classical OOD balanced-accuracy effects. ``Native'' uses the
complete family pools; ``equal count'' block-samples the larger family.
Positive values favour the fidelity-kernel family.}
\label{supp:tab_factorial}
\scriptsize
\setlength{\tabcolsep}{3.5pt}
\begin{tabular}{lrrrr}
\toprule
Regularization and selector & Customary native & Customary equal & Extended native & Extended equal \\
\midrule
$C=1$, OOD oracle & $+0.0131$ & $+0.0094$ & $-0.0010$ & $+0.0021$ \\
$C=1$, ID validation & $+0.0118$ & $+0.0118$ & $-0.0079$ & $-0.0068$ \\
Train-CV, OOD oracle & $+0.0018$ & $-0.0012$ & $-0.0066$ & $-0.0040$ \\
Train-CV, ID validation & $-0.0074$ & $-0.0084$ & $-0.0040$ & $-0.0047$ \\
\bottomrule
\end{tabular}
\end{table}

Supplementary Table~\ref{supp:tab_factorial_contrasts} summarizes paired axis changes and
pairwise differences-in-differences. Each axis row averages its paired change
over the other eight factorial cells; each interaction row averages over the
four cells of the remaining two axes. The relatively large
regularization-by-reference interaction shows why the axis averages are not
interpreted as independent, additive, or causal contributions.

\begin{table}[htbp]
\centering
\caption{Descriptive factorial axis changes and pairwise interactions.
Changes are level-two minus level-one quantum-minus-classical effects, so a
negative value is a shift towards the classical family. Interactions are
differences-in-differences. No population inference, causal attribution, or
familywise decision rule is attached to these finite-design summaries.}
\label{supp:tab_factorial_contrasts}
\scriptsize
\setlength{\tabcolsep}{5pt}
\begin{tabular}{lr}
\toprule
Finite-design summary & Mean change \\
\midrule
Regularization: $C=1\to$ train-CV & $-0.0084$ \\
Selection: OOD oracle $\to$ ID validation & $-0.0037$ \\
Reference: customary $\to$ extended & $-0.0079$ \\
Budget: native $\to$ equal count & $-0.0002$ \\
\addlinespace
Interaction: reference $\times$ budget mode & $+0.0034$ \\
Interaction: regularization $\times$ budget mode & $-0.0007$ \\
Interaction: regularization $\times$ reference & $+0.0139$ \\
Interaction: regularization $\times$ selection & $+0.0000$ \\
Interaction: selection $\times$ budget mode & $+0.0001$ \\
Interaction: selection $\times$ reference & $+0.0004$ \\
\bottomrule
\end{tabular}
\end{table}

\section*{Supplementary Results: Network-field ablation}
\label{supp:shortcut}

The frozen q1000 sensitivity removes port-, protocol-, and service-related
fields before fitting the shared train-only embedding. The ablation changes
absolute task difficulty, especially for the ToN-IoT campaign shift, and can
affect the two kernel families differently. Supplementary Table~\ref{supp:tab_shortcut}
therefore reports original, ablated, and paired-change effects separately.

\begin{table}[htbp]
\centering
\caption{Port/protocol/service-field ablation under train-CV, P1$'$, and
blocked 60-candidate matching (SVC). Scenario-group brackets are pointwise
95\% conditional cluster-$t$ intervals over five q-split clusters and are not
simultaneous. The final-row bracket is the range of the UNSW-NB15 and ToN-IoT
source-dataset changes, not a confidence interval.}
\label{supp:tab_shortcut}
\scriptsize
\setlength{\tabcolsep}{3pt}
\resizebox{\textwidth}{!}{%
\begin{tabular}{lrrr}
\toprule
Scenario & Original effect & Ablated effect & Ablated minus original \\
\midrule
UNSW-DoS (campaign) & $-0.0021\;[-0.0208,+0.0167]$ & $-0.0108\;[-0.0231,+0.0016]$ & $-0.0087\;[-0.0234,+0.0061]$ \\
UNSW-DoS (constructed) & $+0.0014\;[-0.0132,+0.0161]$ & $+0.0008\;[-0.0264,+0.0280]$ & $-0.0006\;[-0.0234,+0.0221]$ \\
UNSW-Recon (campaign) & $-0.0028\;[-0.0275,+0.0220]$ & $+0.0072\;[-0.0091,+0.0234]$ & $+0.0099\;[-0.0268,+0.0466]$ \\
UNSW-Recon (constructed) & $-0.0138\;[-0.0284,+0.0007]$ & $-0.0074\;[-0.0289,+0.0141]$ & $+0.0065\;[-0.0186,+0.0315]$ \\
ToN-IoT (campaign) & $-0.0007\;[-0.0111,+0.0097]$ & $-0.0801\;[-0.1717,+0.0114]$ & $-0.0794\;[-0.1753,+0.0165]$ \\
ToN-IoT (constructed) & $-0.0029\;[-0.0095,+0.0037]$ & $-0.0129\;[-0.0516,+0.0259]$ & $-0.0100\;[-0.0459,+0.0259]$ \\
\midrule
Two-source-dataset-equal & $-0.0031$ & $-0.0245$ & $-0.0215\;[-0.0447,+0.0018]$ \\
\bottomrule
\end{tabular}
}
\end{table}

\clearpage
\section*{Supplementary Results: Circuit resources}
\label{supp:circuits}

\begin{table}[htbp]
\centering
\caption{Complete logical circuit audit. Counts follow Qiskit 2.3.1
decomposition to \texttt{rz}, \texttt{sx}, \texttt{x}, and \texttt{cx} at
optimization level 0 without device routing. ``Map'' is one feature-map state
preparation; ``fidelity'' is its compute--uncompute template for one pair of
inputs. One-qubit counts exclude barriers and measurements. These
compiler-conditional logical counts are not device costs.}
\label{supp:tab_circuits}
\scriptsize
\setlength{\tabcolsep}{3.4pt}
\begin{tabular}{llrrrrrr}
\toprule
Map & Qubits & Map depth & Map 1q & Map CX & Fidelity depth & Fidelity 1q & Fidelity CX \\
\midrule
\multirow{5}{*}{ZZ, 1 rep.}
 & 4  & 19 & 22  & 12  & 38  & 44  & 24 \\
 & 6  & 31 & 39  & 30  & 62  & 78  & 60 \\
 & 8  & 43 & 60  & 56  & 86  & 120 & 112 \\
 & 10 & 55 & 85  & 90  & 110 & 170 & 180 \\
 & 12 & 67 & 114 & 132 & 134 & 228 & 264 \\
\addlinespace
\multirow{5}{*}{ZZ, 2 reps.}
 & 4  & 35  & 44  & 24  & 70  & 88  & 48 \\
 & 6  & 53  & 78  & 60  & 106 & 156 & 120 \\
 & 8  & 71  & 120 & 112 & 142 & 240 & 224 \\
 & 10 & 89  & 170 & 180 & 178 & 340 & 360 \\
 & 12 & 107 & 228 & 264 & 214 & 456 & 528 \\
\addlinespace
\multirow{5}{*}{Pauli-X/Z, 1 rep.}
 & 4  & 11 & 44  & 0 & 22 & 88  & 0 \\
 & 6  & 11 & 66  & 0 & 22 & 132 & 0 \\
 & 8  & 11 & 88  & 0 & 22 & 176 & 0 \\
 & 10 & 11 & 110 & 0 & 22 & 220 & 0 \\
 & 12 & 11 & 132 & 0 & 22 & 264 & 0 \\
\addlinespace
\multirow{5}{*}{Z, 2 reps.}
 & 4  & 8 & 32 & 0 & 16 & 64  & 0 \\
 & 6  & 8 & 48 & 0 & 16 & 96  & 0 \\
 & 8  & 8 & 64 & 0 & 16 & 128 & 0 \\
 & 10 & 8 & 80 & 0 & 16 & 160 & 0 \\
 & 12 & 8 & 96 & 0 & 16 & 192 & 0 \\
\bottomrule
\end{tabular}
\end{table}

\section*{Supplementary Results: External fixed tasks}
\label{supp:external}

All 30 prospectively specified task--size--seed units passed the acquisition,
split-integrity, nesting, and leakage audits. They contain 37,800 split-level
results. Under SVC, size-equal controlled effects for College Scorecard,
diabetes readmission, and ACS income are $-0.0160$, $-0.0066$, and $-0.0131$;
their protocol contractions are $+0.0393$, $+0.0105$, and $+0.0126$.
The task-equal controlled effect is $-0.0119$ (95\% conditional seed-cluster
interval $[-0.0205,-0.0033]$, $n=5$ seeds per task), and contraction is
$+0.0208$ ($[+0.0113,+0.0307]$). Under GPC, the task-equal controlled effect
is $+0.0077$ ($[-0.0015,+0.0183]$) and contraction is $+0.0324$
($[+0.0223,+0.0427]$). These fixed-task results corroborate protocol
sensitivity, not a universal family ordering.

\section*{Supplementary Results: Predictive uncertainty}
\label{supp:uncertainty}

Predictive entropy rises from ID to OOD for both honestly selected families
in every scenario-group. Expected calibration error changes moderately and
comparably (Supplementary Figure~\ref{supp:fig_uncertainty}). Together with the controlled
accuracy results, these diagnostics provide no consistent evidence in the
fixed cases of a calibration benefit unique to the selected fidelity kernels.
Entropy increase alone is not evidence of good calibration.

\begin{figure}[ht]
\centering
\includegraphics[width=\textwidth]{fig_v4_uncertainty.pdf}
\caption{Mean ID-to-OOD change in predictive entropy (left) and 15-bin
expected calibration error (right) for the P1$'$-selected configuration of
each family. Bars aggregate fixed scenario-groups by source dataset and are
descriptive; no population interval is implied.}
\label{supp:fig_uncertainty}
\end{figure}

\section*{Supplementary Results: Repeated finite-shot sensitivity}
\label{supp:shots}

Supplementary Figure~\ref{supp:fig_shots} displays all 2,880 finite-shot evaluations:
eight fixed exact Gram matrices, four shot counts, 30 measurement replicates,
and three PSD-handling conditions. The bands make the conditional Monte Carlo
spread visible rather than treating diversity across data configurations as
measurement replication. The independent-square heuristic and coherent
train-based Nystr{\"o}m extension often narrow or shift the distribution but
do not impose a uniform direction. Exact-statevector OOD accuracy is
approached unevenly because each sampled training matrix reselects a discrete
SVC $C$.

\begin{figure}[ht]
\centering
\includegraphics[width=\textwidth]{fig_v6_shots_mc.pdf}
\caption{Repeated finite-shot OOD change from the exact fixed-configuration
endpoint. Panels show one fixed scenario-group Gram matrix. Lines are medians
over 30 independently seeded measurement replicates and shaded regions are
central 95\% Monte Carlo intervals ($n=30$) conditional on that matrix.
Orange squares use the sampled indefinite matrix; blue circles use independent
PSD projections of train and OOD squares while leaving rectangular blocks
unchanged; green triangles use the coherent train-eigenspace Nystr{\"o}m
extension. The model includes fidelity sampling only, not device noise.}
\label{supp:fig_shots}
\end{figure}

Across the eight group medians, the perturbed/exact effective-rank ratio
declines from 1.93 at 128 shots to 1.13 at 8,192 shots. OOD-KTA under the
Nystr{\"o}m extension is estimator-dependent because its OOD square is derived
from OOD--train features rather than independently sampled; it is not treated
as a mechanism or compared directly with the other KTA estimators. Retained
rank, eigenvalue tolerance, blockwise Frobenius changes, negative-spectrum
fractions, selected $C$, and every replicate endpoint are provided in
\texttt{results/v6/shots\_mc/}.

\begin{table}[ht]
\centering
\caption{Projected sampling resources for the repeated finite-shot
sensitivity. One case--replicate contains 1,624,250 distinct fidelity
parameterizations. ``All eight'' multiplies one shot level by eight fixed
cases and 30 measurement replicates; ``entangling four'' retains only the four
selected ZZ cases. The other four selected configurations are
separable-product maps whose direct classical factorization can bypass circuit
sampling. The three PSD conditions reuse these measurements. Counts exclude
device routing, calibration, mitigation, queueing, and repeated-job overhead;
no hardware execution was performed.}
\label{supp:tab_shot_resources}
\scriptsize
\setlength{\tabcolsep}{4pt}
\begin{tabular}{rrrrr}
\toprule
Shots/fidelity & Shots/case--replicate & Estimation instances/all eight & Shots/all eight & Shots/entangling four \\
\midrule
128   & 207,904,000    & 389,820,000 & 49,896,960,000 & 24,948,480,000 \\
512   & 831,616,000    & 389,820,000 & 199,587,840,000 & 99,793,920,000 \\
2,048 & 3,326,464,000  & 389,820,000 & 798,351,360,000 & 399,175,680,000 \\
8,192 & 13,305,856,000 & 389,820,000 & 3,193,405,440,000 & 1,596,702,720,000 \\
\midrule
All four & 17,671,840,000 & 1,559,280,000 & 4,241,241,600,000 & 2,120,620,800,000 \\
\bottomrule
\end{tabular}
\end{table}

\newpage
\section*{Supplementary Results: Additional geometry and rank-matching displays}

\begin{figure}[ht]
\centering
\includegraphics[width=0.72\textwidth]{fig_v4_rankmatched.pdf}
\caption{Nearest-rank pairing yields negative median quantum-minus-classical
differences in all eight fixed scenario-groups. Each quantum configuration is
paired, with replacement, to the nearest-effective-rank classical
configuration within run and dimension, subject to rank ratio $\leq1.25$.
Annotations give the fraction of retained pairs in which the quantum member
is more accurate. The pairing is observational and does not imply a causal
family effect.}
\label{supp:fig_rankmatched}
\end{figure}

\begin{figure}[ht]
\centering
\includegraphics[width=\textwidth]{fig_v4_mechanism.pdf}
\caption{The geometry--OOD association is regime-dependent. Points are median
within-unit Spearman correlations between each descriptor and OOD balanced
accuracy for SVC. Effective rank is strong in some regimes and weak or
negative in others; OOD alignment is more consistently positive but reuses
the target labels. These descriptive associations do not identify a causal
mechanism.}
\label{supp:fig_mechanism}
\end{figure}

\begin{table}[ht]
\centering
\caption{Zero-label certificate across the frozen nested reference breadths.
The endpoint column is the median across the 16 cell-level medians; the range
is across those cell medians. Final columns count cells whose median upper
endpoint is at or below each materiality threshold. Every 30- and
60-candidate cell uses 5,000 whole-block orderings; breadth 115 contains the
complete family in every ordering.}
\label{supp:tab_v11_breadth}
\scriptsize
\setlength{\tabcolsep}{5pt}
\begin{tabular}{rrrrrrr}
\toprule
$B$ & Blocks & Median $U_B^{(0)}$ & Cell range & $U\leq.005$ & $U\leq.010$ & $U\leq.020$\\
\midrule
30  &  6 & .042 & .004--.094 & 2 & 3 & 3\\
60  & 12 & .034 & .002--.088 & 2 & 4 & 5\\
115 & 23 & .034 & .002--.088 & 2 & 4 & 5\\
\bottomrule
\end{tabular}
\end{table}

\clearpage
\section*{Supplementary Results: Machine-readable result map}
\label{supp:artifact}

The retrospective certificate artifacts and manuscript gates are reproduced by
\begin{verbatim}
python scripts/reproduce_v9.py --stage analysis
python scripts/reproduce_v9.py --stage report
python scripts/reproduce_v9.py --stage audit
python scripts/reproduce_v10.py --stage all
python scripts/reproduce_v11.py --stage all
\end{verbatim}
The exact-prediction export is separately available as
\texttt{python scripts/reproduce\_v9.py --stage export}; it skips complete,
integrity-gated case artifacts unless forced. Target prediction matrices,
separated audit labels, sharp envelopes, evidence-frontier curves, and
finite-shot certificate outputs are under
\texttt{results/v9/partial\_identification/}. Circuit-aware strata and the
evaluation-choice factorial remain under \texttt{results/v8/}; prior inference,
budget, rank, and finite-shot inputs remain under \texttt{results/v6/}.
Frozen earlier artifacts are unchanged. The development repository is
\url{https://github.com/roberto-fernandez-barrios/kernel_shift_framework},
and the immutable v1.1.6 archive is
\url{https://doi.org/10.5281/zenodo.21776862}; the all-version concept DOI is
\url{https://doi.org/10.5281/zenodo.19147649}. The exact v0.8.0 predecessor is
identified by \url{https://doi.org/10.5281/zenodo.21717074}, and the exact
v1.1.5 predecessor is \url{https://doi.org/10.5281/zenodo.21764577}; the exact
v1.1.4 predecessor is \url{https://doi.org/10.5281/zenodo.21759058}. The
prospective replication inputs and outputs are under
\texttt{results/v10/gate2\_prospective/}; the frozen specification,
acquisition audits, prediction locks, opaque label hashes, one-time opening
record, all policy draws, technical-unavailability records, and final outcome
are included in the v1.1.6 archive together with the bounded-loss tests,
consolidation specification, manuscript sources, and final PDFs.

\end{document}
